% Note for D. M. Gordon: a summary of run 1 (tag signed-difference-sets/run1).
% Build: latexmk note.tex (LuaLaTeX, see latexmkrc). Numbers: python numbers.py.
\documentclass[11pt]{amsart}

\usepackage{fontspec}
\setmonofont{DejaVu Sans Mono}[Scale=MatchLowercase]
\usepackage{amssymb}
\usepackage{booktabs}
\usepackage{tabularx}
\usepackage{longtable}
\usepackage{fancyvrb}
\usepackage[margin=2.5cm]{geometry}
\usepackage[hidelinks]{hyperref}

\newcommand{\tableRows}{70{,}543}
\newcommand{\sourceOpen}{67{,}823}
\newcommand{\openAfterHeWu}{22{,}459}
\newcommand{\ourRows}{16{,}101}
\newcommand{\ourNone}{16{,}051}
\newcommand{\ourExists}{50}
\newcommand{\corrections}{11}
\newcommand{\stillOpen}{6{,}347}
\newcommand{\pendingSets}{16}
\newcommand{\settledPercent}{91.0}
\newcommand{\priorPercent}{68.2}
\expandafter\newcommand\csname zrows1\endcsname{1}
\expandafter\newcommand\csname zrows8\endcsname{3{,}858}
\expandafter\newcommand\csname zrows2\endcsname{1}
\expandafter\newcommand\csname zrows17\endcsname{213}
\expandafter\newcommand\csname zrows18\endcsname{297}
\expandafter\newcommand\csname zrows3\endcsname{1}
\expandafter\newcommand\csname zrows4\endcsname{1}
\expandafter\newcommand\csname zrows5\endcsname{1}
\expandafter\newcommand\csname zrows6\endcsname{1}
\expandafter\newcommand\csname zrows33\endcsname{9}
\expandafter\newcommand\csname zrows15\endcsname{2{,}603}
\expandafter\newcommand\csname zrows24\endcsname{1{,}457}
\expandafter\newcommand\csname zrows26\endcsname{2}
\expandafter\newcommand\csname zrows11\endcsname{377}
\expandafter\newcommand\csname zrows12\endcsname{10}
\expandafter\newcommand\csname zrows10\endcsname{1{,}387}
\expandafter\newcommand\csname zrows27\endcsname{1{,}620}
\expandafter\newcommand\csname zrows13\endcsname{3}
\expandafter\newcommand\csname zrows9\endcsname{2{,}213}
\expandafter\newcommand\csname zrows16\endcsname{236}
\expandafter\newcommand\csname zrows29\endcsname{271}
\expandafter\newcommand\csname zrows28\endcsname{716}
\expandafter\newcommand\csname zrows30\endcsname{259}
\expandafter\newcommand\csname zrows23\endcsname{19}
\expandafter\newcommand\csname zrows32\endcsname{7}
\expandafter\newcommand\csname zrows14\endcsname{7}
\expandafter\newcommand\csname zrows20\endcsname{5}
\expandafter\newcommand\csname zrows19\endcsname{144}
\expandafter\newcommand\csname zrows22\endcsname{132}
\expandafter\newcommand\csname zrows25\endcsname{222}
\expandafter\newcommand\csname zrows31\endcsname{26}
\expandafter\newcommand\csname zrows21\endcsname{1}
\expandafter\newcommand\csname zrows7\endcsname{1}
\expandafter\newcommand\csname openBelow50\endcsname{19}
\expandafter\newcommand\csname openBelow100\endcsname{194}
\expandafter\newcommand\csname openBelow200\endcsname{790}
\expandafter\newcommand\csname openBelow500\endcsname{5{,}344}
\expandafter\newcommand\csname openBelowMax\endcsname{0}

\newcommand{\zrows}[1]{\csname zrows#1\endcsname}
\newcommand{\repo}{\url{https://github.com/UltimateSeeSharp/solverz}}
\newcommand{\todo}[1]{\textbf{[TODO: #1]}}

\title[Open cases of the SDS table]{Settling open cases of the signed difference set table,
with Lean-checked proofs: a summary}
\author{\todo{your name}}
\address{Independent researcher}
\email{\todo{email}}
\date{\todo{date}}

\begin{document}

\begin{abstract}
Of the \openAfterHeWu{} cases of Gordon's signed difference set table~\cite{gordon} that
remain open after He and Wu~\cite{he-wu}, we settle \ourRows: \ourNone{} have no signed difference set and \ourExists{} have
one. A further \corrections{} cases listed as open are settled by published constructions.
Every result is a theorem checked by the Lean~4 proof assistant. This note summarizes the
results, the database corrections and how to check them; it is not a paper.
\end{abstract}

\maketitle

\section{How this work was done}

The author is a software developer, not a mathematician. The research was done by AI agents
(Anthropic's Claude models): they searched the literature, proposed and proved the
nonexistence rules, ran the searches and wrote the Lean proofs. Two safeguards decide what
counts. First, nothing is reported unless Lean~4 checks it, with no unproved steps and no
axioms beyond the three standard ones of Lean's logic. Second, before any result was accepted,
a separate agent tried to break it: the proof, the Lean statement, the rows it settles, and
the credit to earlier work. All code, proofs and data are public (\repo).

The author's part was to run and direct this process and to approve each result. Errors in
this note's prose, credits and attributions are possible; the Lean statements are the
authoritative form of every claim.

\section{Results}

We use the notation of~\cite{gordon}: an SDS with parameters $(v,k,\lambda)$ in an abelian
group $G$ of order $v$, $n = k - \lambda$, and $s^2 = k + \lambda(v-1)$ its squared sign sum.
Most rules work with the \emph{coset sums} $a_h = \sum_{\phi(x) = h} D(x)$ of $D$ on a
quotient $\phi \colon G \to H$, with $c = v/|H|$; they satisfy $|a_h| \le c$ and
$\sum_h a_h a_{h+g} = \lambda c$ for $g \neq 0$.

Table~\ref{tab:rules} lists the rules. Each is one Lean theorem that holds for every group;
a row of the table is settled by checking the rule's hypotheses for that row inside Lean.
Rows are counted once, by the first rule that settles them. ``Ours'' means we found no
earlier statement for signed difference sets; it does not mean the idea is new.

\begingroup\small
\begin{longtable}{@{}l p{0.5\textwidth} r p{0.26\textwidth}@{}}
\caption{Rules settling open cases (no SDS).}\label{tab:rules}\\
\toprule
& Rule & Rows & Origin \\ \midrule
\endfirsthead
\toprule
& Rule & Rows & Origin \\ \midrule
\endhead
Z8 & Self-conjugacy: $p^{2j} \mid n$, $p \nmid |H|$, $p$ self-conjugate modulo $\exp H$; the
  $a_h$ agree modulo $p^j$, impossible when $p^j > 2c$ & \zrows{8} & Turyn~\cite{turyn};
  SDS form ours (odd exponent:~\cite{he-wu}) \\
Z9 & Moments: the same congruence with $\sum a_h = \pm s$ and $\sum a_h^2 = k + \lambda(c-1)$
  & \zrows{9} & ours \\
Z10 & Box bound: the same, with $|a_h| \le c$ & \zrows{10} & ours \\
Z22 & The congruence for several self-conjugate primes at once (CRT) & \zrows{22} & ours \\
Z11 & Support parity: $v$ odd, $\lambda$ even, $d(d-1) < v-1$ with $d = v - k$ & \zrows{11}
  & classical for $\lambda = 0$~\cite{wallis-whiteman}; general $\lambda$ ours \\
Z16 & Marked elements per coset (an LP certificate) & \zrows{16} & ours \\
Z17 & Real characters on a quotient $C_2^r$ & \zrows{17} & extends~\cite[Thm.~4.2]{he-wu} \\
Z18 & Characters on a quotient of exponent 4 & \zrows{18} & ours \\
Z19 & Characters on a quotient of exponent 3 & \zrows{19} & the ramified prime above 3 \\
Z25 & Sylow quotient of order $p^t$: $p^{\min(t, \lceil e/2 \rceil)} \mid s$, $e = v_p(n)$;
  for $p = 2$ a sharper equality case & \zrows{25} & ours, on Turyn's argument \\
Z24 & Ma's lemma on a quotient $C_{p^j} \times K$, $p$ self-conjugate on $K$:
  impossible when $p^{\lfloor e/2 \rfloor} > 2c$ & \zrows{24} & Ma's lemma~\cite{pott};
  SDS form ours \\
Z29 & Ma's lemma at its edge ($p = 2$, $4 \mid n$): $8 \mid 2s' - (k + \lambda(2c-1))$;
  for $c = 1$, a halved SDS & \zrows{29} & halving known for weighing matrices~\cite{tan};
  general $c$ ours \\
Z32 & The halved coset sums for any $c$ (even contraction) & \zrows{32} & ours \\
Z15 & Multiplier theorem ($\gcd(v,n) = 1$, $n$ squarefree), then an exhaustive search over
  multiplier orbits & \zrows{15} & McFarland, as in~\cite[Thm.~5.1]{gordon};
  fixed translate~\cite{mcfarland-rice} \\
Z27 & The same for every $n$ (Lean proof through Witt vectors) & \zrows{27} & as Z15 \\
Z28 & Multiplier-fixed coset sums on a quotient, searched exhaustively & \zrows{28} & ours \\
Z31 & Self-conjugate primes enlarge the multiplier group of the coset sums & \zrows{31} &
  for difference sets~\cite{gordon-schmidt,miller}; SDS form ours \\
Z30 & Field descent without ramification: $\varphi(m)\,n \le m F c^2$ on a quotient
  $\mathbb{Z}_m$, $\gcd(m,n) = 1$ & \zrows{30} & Schmidt's field descent~\cite{tan};
  this form ours \\
Z1, Z2, & Exhaustive search and contraction for & 4 & ours \\
Z26 & single small cases & & \\
Z33 & Contraction to a quotient, then a lift coset by coset & \zrows{33} & ours \\
\bottomrule
\end{longtable}
\endgroup

The \ourExists{} cases with a set are listed in Appendix~\ref{app:sets}. They come from
exhaustive and orbit searches, from cyclotomic classes, from partial difference sets
($P - \{0\}$ is an SDS when $\lambda' - \mu = 2$), and from signed combinations of relative
difference sets, among them a cyclic family $(n((n+1)^2+1),\, (n+1)^2+n,\, n)$ for every prime
power $n+1$. The largest is a $(361, 316, 141)$ set in $C_{19}^2$.

\section{Corrections to the database}

These \corrections{} cases are listed as open in \texttt{sds.json} (commit \texttt{e3bf810}),
although published constructions settle them. For each we give a set checked in Lean.

\begin{center}\small
\begin{tabular}{@{}lll@{}}
\toprule
Cases & Parameters & Construction \\ \midrule
N10864, N12056, N13963, & $(211,106,51)$, $(223,112,54)$, $(239,120,58)$, & \cite[Prop.~2.3]{gordon}
  (Paley and zero) \\
N16190, N18527, N19488, & $(251,126,61)$, $(263,132,64)$, $(271,136,66)$, & \\
N49190 & $(419,210,103)$, all cyclic & \\
N67 & $(27,14,5)$ in $C_3^3$ & \cite[Prop.~2.3]{gordon} in $\mathrm{GF}(27)$ \\
N73 & $(27,17,8)$ in $C_3^3$ & \cite[Thm.~4.1]{gordon}, $q = 3$ \\
N49624 & $(421,316,25)$ cyclic & \cite[Prop.~5.5(2)]{he-chen-ge} \\
N14471 & $(243,82,1)$ in $C_3^5$ & \cite[Thm.~4.1]{he-chen-ge} with $x_1 = 0$ \\
\bottomrule
\end{tabular}
\end{center}

For N14471, Example~4.2 of~\cite{he-chen-ge} as printed (with $x_1 = (1,0)$) is not a signed
set; the theorem with $x_1 = 0$ gives one.

\section{How to check it}

The Lean library lives in the repository \repo, in the folder
\path{problems/signed-difference-sets/lean} (Lean~4 v4.35.0 with Mathlib). The statement every
result refers to is:

\begin{Verbatim}[fontsize=\small]
def IsSDS (k : ℕ) (lam : ℤ) (D : G → SignType) : Prop :=
  (Finset.univ.filter fun x => D x ≠ 0).card = k ∧
    ∀ g : G, g ≠ 0 → ∑ x, ((D x : ℤ) * (D (x + g) : ℤ)) = lam

def HasSDS (G : Type*) [AddCommGroup G] [Fintype G] (k : ℕ) (lam : ℤ) : Prop :=
  ∃ D : G → SignType, IsSDS k lam D
\end{Verbatim}

\noindent A single case then reads, for example,

\begin{Verbatim}[fontsize=\small]
theorem not_hasSDS : ¬ HasSDS (ZMod 3 × ZMod 3) 8 1
\end{Verbatim}

\noindent The command \texttt{lake build --wfail} checks everything; \texttt{\#print axioms} on any result
shows only \texttt{propext}, \texttt{Classical.choice} and \texttt{Quot.sound}. No result uses
\texttt{native\_decide}: every computation is checked by Lean's kernel.

\section{What is still open}

\stillOpen{} cases remain: \csname openBelow50\endcsname{} with $v < 50$,
\csname openBelow100\endcsname{} with $50 \le v < 100$, \csname openBelow200\endcsname{} with
$100 \le v < 200$ and \csname openBelow500\endcsname{} with $200 \le v < 500$. Roughly, they
fall into three kinds: cases where a multiplier applies but the orbit search is too large;
cases where a rule above fails by a small margin; and a core where no multiplier applies and
every congruence modulus is below $c$. We expect most of the core to have no SDS, but none of
the local arguments above can show it.

\appendix
\section{The cases settled with a set}\label{app:sets}

\begingroup\small
\begin{longtable}{@{}llll@{}}
\toprule
Case & $(v,k,\lambda)$ & Group & Solution \\ \midrule
\endfirsthead
\toprule
Case & $(v,k,\lambda)$ & Group & Solution \\ \midrule
\endhead
\csname @@input\endcsname sets % TeX's primitive \input: LaTeX's breaks inside a table
\bottomrule
\end{longtable}
\endgroup

\bibliographystyle{amsplain}
\bibliography{refs}

\end{document}
